\documentclass[11pt]{article}
\usepackage[a4paper,margin=28mm]{geometry}
\usepackage{amsmath,amssymb,amsthm,booktabs,hyperref}
\newtheorem{theorem}{Theorem}[section]
\newtheorem{lemma}[theorem]{Lemma}
\newtheorem{proposition}[theorem]{Proposition}
\newtheorem{corollary}[theorem]{Corollary}
\theoremstyle{definition}\newtheorem{definition}[theorem]{Definition}
\newcommand{\dist}{\operatorname{dist}}
\newcommand{\E}{\mathcal E}
\title{Six quartic graphs at every minimum order of diameter $4r$}
\author{Human Compression project: mathematical preparation}
\date{31 October 2026}
\begin{document}\maketitle

\paragraph{Scope and attribution.}
The equality-chain and peeling argument was proposed by the project's
coordinating agent and expanded and checked here by hand. The proof
uses the standard independent-set degree bound, geodesic shortcuts,
graph eccentricities and isomorphism invariants. It makes no literature
novelty claim, uses no known finite census count as a premise, and
does not claim a completed original manuscript, new intake, conjecture
credit, compression measurement or Lean verification. No Source,
graph enumeration, scientific data or TeX build was executed. The
edition is a static mathematical draft awaiting independent review
and adoption. The preceding general minimum-order proof and the
diameter-eight equality proof are retained separately~\cite{general,eight}.
All ingredients needed for the ordinary classification are given below.

Throughout, graphs are finite, simple and connected. A port of a
half means an intrinsic vertex of degree two in that half; every
other half vertex will have degree four. We count ordinary unrooted
graph isomorphism classes, not geodesics, colorings or marked endpoints.

\begin{theorem}\label{main}
For every integer $r\geq2$, the minimum order of a quartic graph of
girth four and diameter $4r$ is $8(r+1)$. There are exactly six ordinary
isomorphism classes at that order. They are obtained from an unordered
pair, with repetition allowed, of three explicitly defined halves of
order $4r+3$, by joining both members of a new two-vertex independent
set to both ports of each half.
\end{theorem}

\section{Equality forces the chain}

Let $\Gamma$ be a triangle-free quartic graph of diameter $4r$,
and choose a geodesic $p_0,p_1,\ldots,p_{4r}$. Set
\[
 z_i=p_{4i},\quad I_i=N(z_i)\quad(0\leq i\leq r),
 \qquad I=\bigcup_{i=0}^r I_i.
\]
Each $I_i$ is independent. Two such neighborhoods are disjoint,
and no edge joins them: otherwise the corresponding geodesic vertices
would be joined by a path of length two or three, shorter than their
distance $4|i-j|$. Thus $I$ is independent of size $4(r+1)$.
The degree sum across an independent set in a regular graph gives
$|V\Gamma|\geq2|I|=8(r+1)$.

Assume equality from now on. The complement $J=V\Gamma\setminus I$
is also independent, because all its degree is used by the edges
from $I$. A vertex in $J$ can meet only one $I_i$ or two consecutive
$I_i$: meeting two nonconsecutive neighborhoods would give a
length-four path between geodesic vertices at distance at least eight.
Let $X_i$ consist of the vertices meeting only $I_i$, and let $F_i$
consist of the vertices meeting both $I_i,I_{i+1}$.

Every vertex of $X_i$ is adjacent to all four vertices of $I_i$.
Thus $X_i$ is a false-twin class containing $z_i$. Put
\[
 x_i=|X_i|\geq1,\qquad q_i=4-x_i.
\]
Connectedness excludes $x_i=4$, since then $X_i\cup I_i$ would
be an isolated $K_{4,4}$. Hence $1\leq q_i\leq3$. Every member
of $I_i$ has all $x_i$ anchor neighbors and exactly $q_i$ other
neighbors in interfaces.

For $1\leq i\leq r-1$, a vertex of $I_i$ cannot meet both
$F_{i-1}$ and $F_i$. If it did, choosing a neighboring vertex in
$I_{i-1}$ through the first interface and in $I_{i+1}$ through the
second would give a path of length six from $z_{i-1}$ to $z_{i+1}$,
whose distance is eight. Therefore
\[
 I_i=I_i^-\sqcup I_i^+,
\]
where all nonanchor neighbors of the first part lie in $F_{i-1}$
and those of the second lie in $F_i$. Both parts are nonempty,
since otherwise the corresponding boundary in the chain of types
would disconnect the graph. All of $I_0$ uses $F_0$ and all of
$I_r$ uses $F_{r-1}$.

\begin{lemma}\label{forced}
For every internal index $1\leq i\leq r-1$,
\[
 |X_i|=2,\qquad q_i=2,\qquad |I_i^-|=|I_i^+|=2.
\]
For $1\leq i\leq r-2$, the interface $F_i$ has two vertices, each
adjacent to all four vertices of $I_i^+\cup I_{i+1}^-$.
\end{lemma}
\begin{proof}
Write $b=|I_1^-|$. Counting all degrees in $F_0$ gives
\[
 4|F_0|=4q_0+bq_1.
\]
Thus $4\mid bq_1$, with $b,q_1\in\{1,2,3\}$. The only possibility
is $b=q_1=2$. Consequently $x_1=2$ and both parts of $I_1$ have size
two.

Suppose the conclusion holds at index $i<r-1$. If
$b'=|I_{i+1}^-|$, the next interface balance is
\[
 4|F_i|=|I_i^+|q_i+b'q_{i+1}=4+b'q_{i+1}.
\]
The same divisibility argument forces $b'=q_{i+1}=2$. Induction proves
the first statement. Then $4|F_i|=4+4$, so $|F_i|=2$.
Each such interface vertex has degree four and all its neighbors
in the four-element set $I_i^+\cup I_{i+1}^-$; its neighborhood is
that entire set. Thus this part of the graph is a $K_{4,2}$.
\end{proof}

The endpoint values $q_0,q_r$ need not equal two. They may be one,
two or three. They are precisely the three endpiece types below.
The internal induction does not restrict them further.

\section{The three base halves}

The left endpiece contains
\[
 X_0\cup I_0\cup F_0\cup I_1^-.
\]
Because $q_1=2$ and $|I_1^-|=2$, its first interface balance gives
$|F_0|=q_0+1$. Its order is therefore
\[
 (4-q_0)+4+(q_0+1)+2=11.
\]
The two vertices of $I_1^-$ are its ports: they have two neighbors
in $F_0$ within the endpiece and two in $X_1$ outside it. Every
other endpiece vertex has degree four internally. The right endpiece
has the same description with reversed indices.

Choose $e=z_0\in X_0$, call $I_0$ the row set, and put
$W=(X_0\setminus\{e\})\cup F_0$. Then $|W|=4$. Every row meets
$e$ and exactly three vertices of $W$. The two port neighborhoods
are two-subsets $P,Q\subseteq W$. If
\[
 \beta(w)=\boldsymbol1_P(w)+\boldsymbol1_Q(w),
\]
the degree-four condition requires a row incidence of $4-\beta(w)$
at $w$. Thus the multiset of omitted vertices of the four rows is
forced: $w$ is omitted exactly $\beta(w)$ times.

\begin{definition}
For $t\in\{0,1,2\}$ choose two two-subsets $P,Q$ of a four-set $W$
with $|P\cap Q|=t$. The base half $H_t$ has one endpoint $e$, four
row vertices, four side vertices $W$, and two ports. The endpoint
meets every row, the ports meet $P,Q$, and each row meets three side
vertices, with omission multiset $\beta$ as above. There are no other
edges. Distinct vertices with the same row omission remain distinct.
\end{definition}

\begin{lemma}\label{base}
These endpieces give exactly the three ordinary unrooted half classes
$H_0,H_1,H_2$. Each is connected and bipartite, has exactly two
degree-two vertices, and all other degrees are four. Every vertex
is within distance three of its port set; the endpoint has distance
exactly three. The number of common neighbors of the two ports is
$t$. In an endpoint presentation the residual parameter is
$q_0=3-t$.
\end{lemma}
\begin{proof}
The preceding derivation proves necessity and the forced row multiset.
Conversely $\sum_w\beta(w)=4$, so four rows with those omissions
exist. Each side degree is $(4-\beta(w))+\beta(w)=4$, each row
degree is $1+3=4$, and each port degree is two. Every side vertex
meets at least two rows; the endpoint joins all rows and each port
has side neighbors, so the half is connected. It is bipartite with
rows and ports in one part and endpoint and side vertices in the other.

Two unordered pairs of two-subsets of a four-set are equivalent
under a side permutation and port exchange exactly when they have
the same intersection size. Indeed the four regions of membership
have sizes $t,2-t,2-t,t$. Such an equivalence extends to an
isomorphism of the forced row multisets. On the other hand the
ports are intrinsic as the only degree-two vertices, and their
common-neighbor number is an ordinary invariant equal to $t$.
Thus all three classes are distinct without marking the endpoint.

Every row, being a three-subset of $W$, meets $P\cup Q$, whose
size is at least two. Hence its distance to a port is at most two.
A supported side vertex has distance one. An unsupported side
vertex has $\beta=0$ and meets all rows, so is within three; the
endpoint is also within three and has no shorter path than
endpoint, row, side, port. Finally the unsupported side vertices
are the $t$ members outside $P\cup Q$. They, together with $e$,
make $X_0$, of size $1+t$. Thus $q_0=4-|X_0|=3-t$.
\end{proof}

\section{Extension and intrinsic peeling}

Let $H$ be a connected half with exactly two degree-two ports
$P=\{u,v\}$, all other degrees four. Define $\E(H)$ as follows.
Add two new vertices $T=\{s,t\}$ and two new ports
$Q=\{u',v'\}$. Join every member of $T$ to every member of
$P\cup Q$, and add no other edge. The old ports now have degree
four, each member of $T$ has degree four, and $Q$ is the new
intrinsic degree-two port set. This operation adds four vertices.

\begin{lemma}\label{peeling}
The extension operation is injective on ordinary unrooted half classes.
More precisely, in $\E(H)$ the two intrinsic ports have the same
two-element neighborhood $T$. Deleting those ports and $T$ recovers
$H$ as a graph, with its two ports again intrinsic by degree.
Any isomorphism of extended halves therefore restricts to an
isomorphism of their predecessors. Conversely every predecessor
isomorphism extends.
\end{lemma}
\begin{proof}
The degree calculation identifies $Q$ uniquely as the degree-two
vertices. Their common neighborhood is exactly $T$, so every graph
isomorphism preserves the four-vertex deletion set. The remaining
graph is precisely $H$; only its old ports lose neighbors, namely
the two members of $T$, and become degree two. This proves the
restriction assertion and injectivity. Since a predecessor
isomorphism preserves its intrinsic port set, it extends by any
bijections on $T$ and $Q$, proving the converse.
\end{proof}

Set $j=r-2$ and write
\[
 H_t^{(r)}=\E^{\,j}(H_t),\qquad t=0,1,2.
\]
These halves have order $11+4j=4r+3$. Their construction is independent
of all matchings or port orderings: every attachment is complete
between the relevant two-element sets.

\begin{lemma}\label{depth}
Each half $H_t^{(r)}$ is connected and bipartite, with exactly two
degree-two ports and all other degrees four. Every vertex is within
distance $2r-1$ of its port set, and its original endpoint has distance
exactly $2r-1$.
\end{lemma}
\begin{proof}
The degree and connectedness claims follow inductively from the
extension definition. The two old ports have the same color in the
base bipartition. Give $T$ the opposite color and $Q$ their original
color at each extension; this preserves bipartiteness.

For an old vertex $z$, every route to the new ports $Q$ must first
reach an old port and then traverse a member of $T$. Hence
\[
 \dist_{\E(H)}(z,Q)=\dist_H(z,P)+2.
\]
The new $T$ vertices have distance one to $Q$ and $Q$ itself has
distance zero. Lemma~\ref{base} starts with maximum distance three,
attained by the endpoint. After $j=r-2$ extensions this gives
$3+2j=2r-1$, with the endpoint still attaining it.
\end{proof}

\section{The full graph and its center}

For $t,u\in\{0,1,2\}$ take disjoint halves
$H_t^{(r)},H_u^{(r)}$. Add two new nonadjacent vertices
$X=\{x,x'\}$, both adjacent to all four ports of the two halves.
Call the resulting graph $G_{t,u}^{(r)}$.

\begin{lemma}\label{metric}
$G_{t,u}^{(r)}$ is simple, connected and quartic, of order $8(r+1)$,
girth four and diameter $4r$. Its center is exactly $X$. Removing
$X$ gives two connected components, each of order $4r+3$, with
intrinsic degree-two port sets.
\end{lemma}
\begin{proof}
Each old port gains its two missing neighbors, each new vertex
has the four ports as neighbors, and all other degrees are unchanged.
The order is $2(4r+3)+2=8(r+1)$. Connectedness and simplicity are
immediate from the construction. The two halves are bipartite,
their ports have the same color, and $X$ is assigned the opposite
color, so the graph is bipartite. The two members of $X$ and any
two ports give a four-cycle. Thus the girth is exactly four.

By Lemma~\ref{depth}, every half vertex has distance at most $2r$
from either member of $X$, and each original endpoint has distance
exactly $2r$ from both. The two members of $X$ have distance two.
Thus every pair of vertices has distance at most $4r$ through $X$.
The two original endpoints lie in different components after
deleting $X$, so their distance is exactly $4r$. This gives the
diameter and eccentricity $2r$ for each member of $X$.

If $z\notin X$ lies in one half, choose the original endpoint
$e$ of the other half. Every $z$--$e$ path meets $X$. The two
members of $X$ are false twins, so their distances to $z$ agree;
their distances to $e$ both equal $2r$. Therefore
\[
 \dist(z,e)=\dist(z,X)+2r>2r.
\]
No such $z$ is central. Hence the center is precisely $X$, not
just a chosen midpoint class. It is a full false-twin class as well:
any further false twin would have the same eccentricity and would
also be central. The component and port statements
follow from construction and Lemma~\ref{depth}.
\end{proof}

\section{Completeness and ordinary identification}

\begin{lemma}\label{complete}
Every triangle-free quartic graph of order $8(r+1)$ and diameter
$4r$ is isomorphic to some $G_{t,u}^{(r)}$. Relative to any chosen
length-$4r$ geodesic, its center is $X_{r/2}$ if $r$ is even and
$F_{(r-1)/2}$ if $r$ is odd.
\end{lemma}
\begin{proof}
The equality structure and Lemma~\ref{forced} leave an eleven-vertex
endpiece of type $H_t$ at each end. Between them the two-element
sets have the chain pattern
\[
 I_1^-\;--\;X_1\;--\;I_1^+\;--\;F_1\;--\;
 I_2^-\;--\;X_2\;--\;\cdots\;--\;X_{r-1}\;--\;I_{r-1}^+,
\]
where each displayed link is a complete bipartite incidence between
two two-element sets. For $r=2$ the chain consists only of
$I_1^-,X_1,I_1^+$. There are no omitted edges: every port and
anchor degree is already four, every interface is the complete
$K_{4,2}$ proved above, and the type decomposition accounts for
every vertex.

Starting from a base endpiece, the addition of each next two-element
anchor or interface class and the next two-element port set is
exactly one application of $\E$. The middle twin class lies at
geodesic position $2r$. The original anchors $X_i$ have position
$4i$ and interfaces $F_i$ position $4i+2$. Thus the middle class
is the one stated, in either parity. There are precisely $r-2$
extension steps from either base endpiece to the ports adjacent
to this middle class. Equivalently their port positions advance
from $3$ to $2r-1$ by two per step. Deleting the middle class
therefore gives $\E^{r-2}(H_t)$ and $\E^{r-2}(H_u)$.
The graph is $G_{t,u}^{(r)}$, and Lemma~\ref{metric} proves that
its middle class is its intrinsic center.
\end{proof}

The construction justifies the positions used in the proof: a path
from the left endpoint first reaches $I_1^-$ at length three and
$X_1$ at length four; each complete link advances one edge. There
is no bypass, since deleting any internal $X_i$ or $F_i$ separates
the two ends. The displayed middle position is therefore a metric
statement, not a choice of names based on a census result.

\begin{proposition}\label{ordinary}
For fixed $r\geq2$,
\[
 G_{t,u}^{(r)}\cong G_{t',u'}^{(r)}
 \quad\Longleftrightarrow\quad
 \{t,u\}=\{t',u'\}\text{ as multisets}.
\]
\end{proposition}
\begin{proof}
An ordinary graph isomorphism preserves all eccentricities, hence
the center $X$. It therefore induces an unordered pair of component
isomorphisms after deleting $X$. The two components have equal
order $4r+3$ and intrinsic degree-two ports. If $r>2$, their ports
have the same two-element neighborhood. Delete those ports and
that neighborhood and repeat this intrinsic operation exactly
$r-2$ times. Lemma~\ref{peeling} shows that an isomorphism restricts
at each step, ending at an isomorphism of the base halves. By
Lemma~\ref{base}, their values of $t$ agree. Thus the unordered
pair of base types is invariant under full ordinary isomorphism.

The stopping rule is part of this argument: stop when the half has
order eleven, after precisely $r-2$ deletions. It is determined by
the original graph order, and is preserved by isomorphisms. One
does not keep deleting merely because a base half happens to have
equal port neighborhoods; in particular $H_2$ does have that
property. No maximal-peeling assumption is made.

Conversely, matching base types gives isomorphisms of the bases,
which extend through all applications of $\E$. A possible exchange
of the two halves is harmless. Both ports in each component are
joined to both members of $X$, so there is no further matching or
alignment parameter. The component isomorphisms extend to a full
graph isomorphism. This proves the equivalence.
\end{proof}

\begin{proof}[Proof of Theorem~\ref{main}]
The independent-neighborhood bound gives order at least $8(r+1)$.
Lemma~\ref{metric} constructs graphs of precisely that order, with
girth four and diameter $4r$. Lemma~\ref{complete} proves exhaustion
at equality. Proposition~\ref{ordinary} gives one class for each
unordered pair of the three base types, with repetition allowed.
Their number is $\binom{3+2-1}{2}=6$.
\end{proof}

\paragraph{Consequences and limits.}
The theorem classifies the quartic equality domain for every
diameter $4r$, rather than extrapolating a finite count at diameter
eight. Its reusable mathematical content is the forced chain,
three base halves, intrinsic center and reversible extension.
This does not enumerate graphs of larger order, other degrees,
or other diameter congruence classes. Computational validation,
formal Lean verification and any comparison of complete source
manuscripts remain separate tasks.

\begin{thebibliography}{9}
\bibitem{general} Project static proof, \emph{Regular graphs of girth
four and diameter divisible by four}, 19,767 bytes, SHA-256
\texttt{782f154caa3d19fe07c70e41856452a1e6032cd702aaf56952acbb7ea40b90c4}.
The general minimum theorem is not used here to assume an ordinary
class count; its equality setup is reproved above.
\bibitem{eight} Project static proof, \emph{A regular-graph bound at
diameter eight and its degree-four equality case}, 13,975 bytes,
SHA-256 \texttt{26f46bb39cb0f22de00b3e7566096289ab18f645c90b3060130c14b3df6fa65e}.
Its three base-half types are proved again here for completeness.
\end{thebibliography}
\end{document}
